\documentclass[journal]{IEEEtran}
\usepackage{amsmath,amssymb,booktabs,bm,cite}

\title{Symmetry-Protected Modal Controllability for Broadband X-Band Magic-T Power Combining}
\author{Author~Names~Omitted~for~Draft}

\begin{document}
\maketitle

\begin{abstract}
A broadband waveguide Magic-T is studied as a modal-controllability problem rather than a scalar S-parameter optimization problem.
Mirror symmetry decomposes the collinear excitation into even and odd parity channels that ideally couple to the H and E arms, respectively.
A full-band real Jacobian is constructed from the complex modal reflections over 9--11.5~GHz, and singular-value analysis is used to quantify independent geometric control directions.
An earlier parameterization is practically ill conditioned with a reported step-normalized condition number of approximately 422.
A later three-parameter central-difference data set at a legacy stepped-post center gives a full-band condition number of approximately 6.95, but the odd/E block remains poorly conditioned at approximately 74.4.
A strict four-parameter identification at the final design seed is therefore proposed by reopening the upper-post radius.
The framework separates modal rank, created by spatially distinct junction perturbations, from spectral rank, created by throat and distributed ridge-transformer sections.
\end{abstract}

\begin{IEEEkeywords}
Magic-T, waveguide combiner, even-odd modes, S-parameter sensitivity, Jacobian conditioning, broadband matching.
\end{IEEEkeywords}

\section{Introduction}
Waveguide Magic-T junctions are attractive for high-power microwave combining because of low conductor loss, high power-handling capability, and their natural sum/difference operation.
Their main difficulty is broadband matching.
Published wideband designs therefore combine several distinct matching elements such as posts, irises, steps, septa, and ridge transitions \cite{peng2024ridge,leal2013compact,ggw2022}.

A large number of geometric variables does not necessarily provide a useful optimization space.
If several parameters generate almost the same electromagnetic response, the local design problem is practically rank deficient.
This behavior was observed in an earlier version of the present X-band combiner, whose step-normalized local Jacobian had condition number approximately 422.

This paper instead works in the natural parity basis and asks whether geometry parameters span independent directions in the two modal reflection trajectories across the complete band.
The intended contributions are: 1) symmetry-based parity decomposition; 2) a full-band Jacobian/SVD controllability measure; 3) a stepped-post modal-rank interpretation; and 4) separation of local modal rank from distributed spectral rank.

\section{Parity Decomposition}
Let ports 1 and 2 be the collinear arms, port 3 the H arm, and port 4 the E arm.
Define
\begin{equation}
c_+=\frac{P_1+P_2}{\sqrt2},\qquad
c_-=\frac{P_1-P_2}{\sqrt2}.
\end{equation}
Let \(\mathcal P\) exchange the two collinear arms.
For a mirror-symmetric geometry,
\begin{equation}
[\mathcal L,\mathcal P]=0,\qquad [S,\mathcal P]=0.
\end{equation}
Hence the ideal modal scattering matrix decomposes as
\begin{equation}
S_m \sim (c_+,H)\oplus(c_-,E).
\label{eq:block}
\end{equation}

Define
\begin{equation}
\Gamma_+(f)=S_{m,11}(f),\qquad
\Gamma_-(f)=S_{m,22}(f).
\end{equation}
For the symmetric collinear section,
\begin{equation}
S_{11}=\frac{\Gamma_++\Gamma_-}{2},\qquad
S_{12}=\frac{\Gamma_+-\Gamma_-}{2},
\end{equation}
so
\begin{equation}
|S_{11}|^2+|S_{12}|^2=
\frac{|\Gamma_+|^2+|\Gamma_-|^2}{2}.
\end{equation}
Thus a small physical \(S_{11}\) can be produced by cancellation and does not prove that both parity channels are individually matched.

\section{Full-Band Modal Controllability}
For geometry parameters \(\bm p=(p_1,\ldots,p_M)^T\), stack the two modal reflections across \(N\) frequencies:
\begin{equation}
\bm F(\bm p)=
[\Re\bm\Gamma_+;\Im\bm\Gamma_+;
 \Re\bm\Gamma_-;\Im\bm\Gamma_-].
\end{equation}
The full-band Jacobian is
\begin{equation}
J=\frac{\partial \bm F}{\partial \bm p}
\in\mathbb R^{4N\times M}.
\end{equation}
The parity blocks are denoted \(J_+\) and \(J_-\).
The local conditioning is measured by
\begin{equation}
\kappa(J)=\frac{\sigma_{\max}(J)}{\sigma_{\min}(J)}.
\end{equation}
Because the parameters have different natural scales, the main comparison uses
\begin{equation}
\widetilde J=JD,
\qquad
D=\mathrm{diag}(\Delta p_1,\ldots,\Delta p_M),
\end{equation}
where \(\Delta p_i\) is the finite-difference half step.

For parameter \(p_i\), define modal selectivity
\begin{equation}
\rho_i^{(+)}=
\frac{\|J_{+,i}\|_2}{\|J_{-,i}\|_2},
\qquad
\rho_i^{(-)}=
\frac{\|J_{-,i}\|_2}{\|J_{+,i}\|_2}.
\end{equation}
Pairwise acute angles between normalized Jacobian columns are used to identify redundant directions.

\section{Stepped-Post Interpretation}
Let \(j_L\) and \(j_U\) be response columns for lower and upper post-radius perturbations.
A uniform through-post enforces
\begin{equation}
\delta r_L=\delta r_U=\delta r,
\end{equation}
hence
\begin{equation}
\delta\bm F=(j_L+j_U)\delta r,
\end{equation}
which provides only one local response direction.

A stepped post instead gives
\begin{equation}
\delta\bm F=j_L\delta r_L+j_U\delta r_U.
\label{eq:step}
\end{equation}
If \(j_L\) and \(j_U\) are not collinear, the reachable response space expands.

This has a field-theoretic interpretation.
For a small normal boundary perturbation \(V_n\), an S-parameter variation has an adjoint/Hadamard boundary-sensitivity form \cite{zhu2009,dadash2014},
\begin{equation}
\delta S=
\int_{\partial\Omega}
G_S(\bm r,\omega)V_n(\bm r)\,dS.
\end{equation}
The lower and upper post segments occupy different field regions, and the even/H and odd/E fields have different spatial distributions.
Their overlap integrals therefore need not be proportional.

\section{Modal Rank and Spectral Rank}
Independent control of \(\Gamma_+\) and \(\Gamma_-\) at one frequency does not guarantee broadband matching.
Around \(f_0\),
\begin{equation}
\Gamma_\pm(f)=
\Gamma_\pm(f_0)+
\Gamma_\pm'(f_0)(f-f_0)+
\frac12\Gamma_\pm''(f_0)(f-f_0)^2+\cdots.
\end{equation}
A local post is predominantly a local reactive or resonant perturbation.
By contrast, a throat or distributed ridge transformer changes both impedance and accumulated propagation phase.
A first-order distributed-reflection picture is
\begin{equation}
\Gamma(f)\approx
\frac12\int
\frac{d\ln Z(z)}{dz}
\exp\left[
-2j\int_0^z\beta(\xi,f)d\xi
\right]dz.
\end{equation}
This motivates
\begin{equation}
\boxed{
\text{stepped local geometry}\rightarrow\text{modal rank},
\quad
\text{throat/ridge}\rightarrow\text{spectral rank}.
}
\end{equation}

\section{Existing Evidence}
The old R11 identification produced step-normalized singular values approximately
\begin{equation}
[0.1646,0.07230,0.002146,0.000390],
\end{equation}
with \(\kappa\approx422\).

The repository also contains a strict three-parameter central-difference set centered at
\begin{equation}
(r_L,h_s,r_U,H)=
(2.0,10.16,0.9,16.25)~\mathrm{mm}.
\label{eq:legacy}
\end{equation}
This is a legacy local-screen center, not the final v7 seed.
At this point,
\begin{equation}
\sigma(J)\approx[4.3147,2.6554,0.6205],
\end{equation}
so
\begin{equation}
\kappa(J)\approx6.95,
\quad
\kappa(J_+)\approx7.06,
\quad
\kappa(J_-)\approx74.4.
\end{equation}
The total-height direction is strongly even-selective there, with
\begin{equation}
\frac{\|J_{+,H}\|_2}{\|J_{-,H}\|_2}\approx47.
\end{equation}

The final v7 single-cell PEC design has H/sum efficiency 0.7801 minimum and 0.8699 mean, and E/difference efficiency 0.7038 minimum and 0.7746 mean.
The complete four-way PEC model has combining efficiency 50.17\% minimum and 78.34\% mean, minimum active return loss about 2.95~dB, and minimum input isolation about 6.94~dB.
The network-to-full-3-D RMS efficiency difference is approximately 1.22 percentage points.
These results suggest that the tree topology is modeled consistently while broadband cell matching remains the dominant limitation.

\section{Strict Final-Seed Identification}
The final v7 seed is
\begin{equation}
(r_L,h_s,r_U,H)=
(1.4,6.2,0.9,16.25)~\mathrm{mm}.
\label{eq:final}
\end{equation}
Existing simulations already provide the center, \(r_L=1.6\)~mm, and \(h_s=7.2\)~mm.
Six additional solves complete a strict four-parameter central difference:
\begin{equation}
r_L=1.2,quad
h_s=5.2,quad
r_U=0.7,1.1,quad
H=15.50,17.00~\mathrm{mm}.
\end{equation}

The decisive test is whether the upper-radius column adds a non-redundant odd/E direction, increases \(\sigma_{\min}(J_-)\), and reduces \(\kappa(J_-)\) over the entire band.

\begin{table}[t]
\centering
\caption{Modal controllability comparison.}
\begin{tabular}{lccc}
\toprule
Case & Full \(\kappa\) & Even/H & Odd/E\\
\midrule
Old R11 & \(\sim422\) & -- & --\\
Legacy center & 6.95 & 7.06 & 74.4\\
Final seed v8 & \textbf{TODO} & \textbf{TODO} & \textbf{TODO}\\
\bottomrule
\end{tabular}
\end{table}

\section{Four-Way Tree}
For four coherent equal inputs,
\begin{equation}
\bm a_\Sigma=\frac12[1,1,1,1]^T,
\end{equation}
the desired state is the global sum mode.
Pairwise difference channels ideally vanish at the first stage, while pairwise sum states combine at the second stage.
This hierarchical sum/difference interpretation explains why coherent amplitude and phase balance can already be accurate even while cell return loss is poor.

\section{Discussion and Conclusion}
The design problem is governed by the rank and conditioning of electromagnetic response directions, not by the raw number of geometric variables.
Mirror symmetry protects the parity decomposition.
The stepped post is useful only if its separated geometric regions create independent modal sensitivities.
The throat and ridge sections provide a complementary frequency-dependent basis required for broadband flattening.

The final v8 identification must be completed before claiming final-seed controllability improvement.
The present simulations are PEC and do not yet establish conductor loss, breakdown, multipactor, thermal margin, or 30-kW peak-power handling.

\begin{thebibliography}{99}
\bibitem{peng2024ridge}
W. Peng, C. Li, H. Wang, and L. Ye,
"An 18--40 GHz ridge waveguide Magic-T using stepped conducting T-junction transition,"
\emph{Electronics}, vol. 13, no. 12, Art. no. 2407, 2024,
doi: 10.3390/electronics13122407.

\bibitem{leal2013compact}
C. A. Leal-Sevillano, J. A. Ruiz-Cruz, J. R. Montejo-Garai, and J. M. Rebollar,
"Compact broadband couplers based on the waveguide magic-T junction,"
in \emph{Proc. European Microwave Conference}, 2013.

\bibitem{ggw2022}
"An H-plane groove gap waveguide Magic-T for X-band applications,"
\emph{Electronics}, vol. 11, no. 24, Art. no. 4075, 2022,
doi: 10.3390/electronics11244075.

\bibitem{zhu2009}
X. Zhu and N. K. Nikolova,
"Accuracy improvement of the S-parameter adjoint sensitivity analysis for shape parameters,"
in \emph{Proc. IEEE MTT-S International Microwave Symposium}, 2009, pp. 529--532,
doi: 10.1109/MWSYM.2009.5165750.

\bibitem{dadash2014}
M. S. Dadash and N. K. Nikolova,
"Analytical S-parameter sensitivity formula for the shape parameters of dielectric objects,"
\emph{IEEE Microwave Wireless Components Letters}, vol. 24, no. 5, pp. 291--293, 2014,
doi: 10.1109/LMWC.2014.2306891.
\end{thebibliography}

\end{document}
